\section*{Chapitre \ref{ch:pl:cloud-against-on-premises} --- Cloud contre infrastructure sur site}

\begin{solution}{exo:pl:cloud-against-on-premises:1}
Le prix horaire de l'instance divisé par ses 96 threads~: 4.46~centimes à la demande, 1.96~centimes réservé pour trois ans sans paiement initial (upfront), et
1.81~centimes en spot (au 28 septembre 2026).
\end{solution}

\begin{solution}{exo:pl:cloud-against-on-premises:2}
Au-delà de 1.886/4.284, soit 44~\%~: en dessous, payer chaque heure une \omterm{def:pl:cloud-against-on-premises:capacity}{capacité réservée} coûte plus cher que payer à la demande les seules heures utilisées.
\end{solution}

\begin{solution}{exo:pl:cloud-against-on-premises:3}
29\,491~\$~: 10~To à 9~centimes par gigaoctet, 40~To à 8.5, 100~To à 7 et les 350~To restants à 5 (avec 1\,024~Go pour 1~To). Le transfert entrant
était gratuit.
\end{solution}

\begin{solution}{exo:pl:cloud-against-on-premises:4}
Les deux sont payées chaque heure, utilisées ou non~; pour toute unité de capacité, c'est donc le prix horaire le plus bas qui l'emporte à tout taux d'utilisation~: 0.40~centime en propriété
contre 1.96~centimes en réservé. La \omterm{def:pl:cloud-against-on-premises:capacity}{capacité réservée} ne l'emporte que lorsque la propriété n'est pas possible --- pas de centre de données, pas le temps d'acheter, un besoin qui dure moins que
la durée de vie du serveur.
\end{solution}

\begin{solution}{exo:pl:cloud-against-on-premises:5}
La durée du balayage est fixée par ses backtests les plus longs, 7.5 heures chacun. Une interruption relance l'un d'eux depuis le début~; si elle survient six heures
après le départ, ce backtest se termine après 13.5 heures, et le balayage avec lui~: 14.4 heures avec la bande du conseiller.
\end{solution}

\begin{solution}{exo:pl:cloud-against-on-premises:6}
Le coût par unité de travail du serveur en propriété est divisé par deux, et le seuil de rentabilité aussi~: 4.5~\% par rapport à la demande, 11~\% par rapport au spot. La vitesse compte autant
que le prix, ce qui explique qu'il faille comparer en backtests par dollar, mesurés.
\end{solution}

\begin{solution}{exo:pl:cloud-against-on-premises:7}
Le balayage ajoute 2\,400 threads au batch de chaque nuit de semaine, un niveau utilisé 24~\% des heures~; le planificateur possède donc 7\,900 threads au lieu de 5\,500, et le coût
annuel s'élève à 291\,003~\$ avec des débordements à la demande, contre 361\,215~\$~: déplacer les travaux flexibles dans les heures que la capacité possédée couvre déjà augmente son taux d'utilisation
et fait plus d'économies que toute négociation de prix.
\end{solution}

\begin{solution}{exo:pl:cloud-against-on-premises:8}
La \omterm{def:pl:cloud-against-on-premises:capacity}{capacité spot} peut être reprise à tout moment~: le travail interactif et tout ce qui n'a pas de points de reprise perd sa progression, et le balayage du chapitre
sans points de reprise a pris deux fois plus de temps à la bande du conseiller. Quant à la charge de base, utilisée en permanence, la capacité en propriété coûte plus de quatre fois moins cher par
heure-thread que le spot, aux prix du modèle.
\end{solution}

\begin{solution}{pb:pl:cloud-against-on-premises:1}
\begin{enumerate}
\item 4.46~centimes à la demande, 2.95 réservé un an, 1.96 réservé trois ans, 1.81 spot~: la grille tarifaire publiée par le fournisseur le 25 septembre
2026 et son flux spot du 28 septembre, divisés par 96 threads.
\item Cités~: le prix catalogue du processeur et le prix moyen de l'électricité commerciale. Supposés~: le reste du serveur (autant que les processeurs),
une durée de vie de quatre ans, 1.3~kW, un coefficient de surcoût de 1.4, 300~\$ d'espace et un deux-centième d'ingénieur par mois.
\item 1\,508~\$ par mois, soit 0.40~centime par heure-thread entièrement utilisée.
\item La capacité fixe (en propriété, réservée) coûte son prix horaire divisé par le taux d'utilisation~; la capacité à l'usage (à la demande, spot) coûte son prix par
heure utilisée.
\item 9.0~\% par rapport à la demande et 22.3~\% par rapport au spot~; 14.5~\% et 35.8~\% avec un serveur à deux fois le prix, 20.0~\% et 49.2~\% à trois fois.
\item Base de 1\,500 threads, 1\,000 de plus aux heures de semaine, 4\,000 de plus les nuits de semaine, 12\,000 de plus pendant huit heures le samedi, vingt rafales
de 8\,000 pendant quatre heures~; moyenne 3\,364, pic 13\,500.
\item À partir de la courbe de durée de charge~: la $k$-ième unité est utilisée sur la part $f(k)$ des heures, et elle est possédée si la propriété coûte moins que $f(k)$ fois
le prix horaire~; réservée selon le même test lorsque la propriété est exclue, sinon achetée à l'heure.
\item 5\,500 threads possédés~; 496 heures ont une demande supérieure à ce niveau.
\item 1\,315\,233~\$ à la demande seule, 940\,460~\$ réservé et à la demande, 361\,215~\$ en propriété et à la demande, 281\,883~\$ en propriété avec des débordements spot.
\item Parce que la demande comporte une base utilisée la plupart du temps et des pics utilisés rarement, de part et d'autre du seuil de rentabilité.
\item Économies de 65~\% par rapport à la demande sur les 30 derniers jours, et fréquence d'interruption dans la bande la plus haute, au-dessus de 20~\% sur le mois
glissant.
\item Avec la bande du conseiller, 94~\$ et 14.4 heures~; à 0.2 par heure, 185~\$ et 59 heures~; contre 203~\$ et 7.5 heures à la demande.
\item La perte est bornée à un segment~: 80--82~\$ et moins de 9 heures à chaque taux.
\item 11\,776~\$ par mois en stockage d'objets standard~; 29\,491~\$ pour sortir une copie une seule fois.
\item Dans l'UE, selon les orientations de l'ABE sur l'externalisation, entre autres une stratégie de sortie documentée pour une fonction critique ou importante.
\item \textbf{Résultat principal.} Aux prix cités et avec le serveur en propriété du modèle, la propriété coûte moins cher que la location à la demande au-dessus de 9~\% d'utilisation
(22~\% par rapport au spot~; 20~\% et 49~\% même à trois fois le prix du serveur). Le balayage du chapitre~13 coûte 203~\$ à la demande et 82~\$ en spot sans
interruption~; à la bande d'interruption la plus haute du fournisseur, il coûte 94~\$ et prend 14.4 heures au lieu de 7.5, et avec des points de reprise toutes les 15 minutes,
80~\$ en moins de 9 heures.
\item La vitesse relative des deux machines sur les backtests propres au cabinet, mesurée~; puis le prix complet du serveur, à partir d'une cotation.
\item Là où s'exécute la plus grande partie du calcul qui le lit~; avec une copie unique faisant foi, et le calcul envoyé vers elle.
\item Une demande dominée par de rares pics importants, pas de centre de données ni de personnel, ou des données déjà dans le cloud.
\item La forme de la demande face au prix de la capacité fixe et de la capacité à l'usage, l'emplacement des données, et la manière dont le travail peut être interrompu.
\end{enumerate}
\end{solution}

\begin{solution}{iq:pl:cloud-against-on-premises:1}
À la demande~: sans engagement, prix le plus élevé, pour un travail irrégulier. Réservé~: engagement d'un ou trois ans à un prix plus bas, pour une charge régulière. Spot~:
capacité excédentaire au prix le plus bas, pouvant être reprise, pour un travail batch interruptible avec points de reprise.
\iqlookfor{Engagement, prix, interruptibilité.}
\end{solution}

\begin{solution}{iq:pl:cloud-against-on-premises:2}
Construire le profil de demande et sa courbe de durée de charge, chiffrer la capacité en propriété par heure en explicitant ses hypothèses, et acheter chaque niveau de capacité
de la façon la moins chère compte tenu de la part des heures où il est utilisé~; vérifier la localisation des données et l'interruptibilité~; comparer en travail par dollar, mesuré.
\iqlookfor{Taux d'utilisation et seuil de rentabilité, pas les prix catalogue.}
\end{solution}

\begin{solution}{iq:pl:cloud-against-on-premises:3}
Rendre chaque tâche relançable~: prévoir des points de reprise pour les tâches longues à intervalles réguliers, garder les tâches idempotentes en écrivant les résultats de façon atomique, laisser
l'ordonnanceur resoumettre le travail interrompu, et répartir sur plusieurs types d'instances et zones.
\iqlookfor{Points de reprise et tâches idempotentes.}
\end{solution}

\begin{solution}{iq:pl:cloud-against-on-premises:4}
Le transférer vers le cloud est gratuit, mais il doit alors y être stocké (environ 12\,000~\$ par mois en stockage standard) et en sortir coûte environ 29\,000~\$ par
copie~; décider où réside la copie faisant foi et envoyer le calcul vers elle, ou ne copier que le sous-ensemble dont la recherche a besoin.
\iqlookfor{\omterm{def:pl:cloud-against-on-premises:gravity}{Gravité des données} et frais de transfert de données.}
\end{solution}

\begin{solution}{iq:pl:cloud-against-on-premises:5}
Un cluster en propriété dimensionné pour le niveau utilisé au-delà de la part de seuil de rentabilité en heures~; un seul ordonnanceur sur les nœuds possédés et ceux du cloud~; des débordements sur spot pour
le travail avec points de reprise et à la demande pour le reste~; les données près du cluster en propriété, avec des sous-ensembles en cache dans le cloud~; l'\omterm{def:pl:cloud-against-on-premises:iac}{infrastructure as code}~; un plan de sortie.
\iqlookfor{Dimensionner la propriété pour la base, déborder vers le cloud pour les pics.}
\end{solution}
